% answers_short.tex —— 2008 年·普通高等学校招生考试·安徽卷·理科 · 简短答案速查
% 由 tools/make_short_answers.py 自 answers.tex 的 \gksolhead/\gkanswer 提取，
% 并与 2008/2008_数学真题_安徽理_答案要点.md 逐题比对后生成，勿手改（改源 MD 后重跑）。

\begin{gknotice}
\textbf{使用说明}：本节只给最终答案便于快速核对，完整推导见解析版。答案为编写者独立推导，\textbf{非官方标准答案}。
\end{gknotice}

\gkpart{选择题}
\begin{center}\begin{tabular}{ccccccccccccc}
\toprule
题号 & 1 & 2 & 3 & 4 & 5 & 6 & 7 & 8 & 9 & 10 & 11 & 12 \\
\midrule
答案 & A & D & B & D & C & A & B & C & B & A & D & C \\
\bottomrule
\end{tabular}\end{center}

\gkpart{填空题}
\begin{itemize}[leftmargin=2.6em]
  \item[\textbf{13.}] $[3,+\infty)$。
  \item[\textbf{14.}] $1$。
  \item[\textbf{15.}] $\dfrac74$。
  \item[\textbf{16.}] $\dfrac{4\pi}3$。
\end{itemize}

\gkpart{解答题}
\begin{itemize}[leftmargin=2.6em]
  \item[\textbf{17.}] （1）$T=\pi$，对称轴 $x=\dfrac{\pi}3+\dfrac{k\pi}2$（$k\in\mathbf Z$）；（2）$\left[-\dfrac{\sqrt3}2,1\right]$。
  \item[\textbf{18.}] （1）证明见解析；（2）$60^\circ$；（3）$\dfrac23$。
  \item[\textbf{19.}] （1）$n=6$，$p=\dfrac12$，$P(\xi=k)=\dfrac{\mathrm C_6^k}{64}$（$k=0,\cdots,6$）；（2）$\dfrac{21}{32}$。
  \item[\textbf{20.}] （1）增区间 $\left(0,\dfrac1{\mathrm e}\right)$，减区间 $\left(\dfrac1{\mathrm e},1\right)$、$(1,+\infty)$；（2）$a>-\mathrm e\ln2$。
  \item[\textbf{21.}] （1）（2）（3）证明见解析。
  \item[\textbf{22.}] （1）$\dfrac{x^2}4+\dfrac{y^2}2=1$；（2）点 $Q$ 恒在定直线 $2x+y-2=0$ 上。
\end{itemize}
